\chapter{Subject-Wise Engineering Concepts}\label{ch:subjects}
This project was meant to bring five courses of the programme to bear on a single
system. This chapter shows, course by course, which concepts were used, where in the
project they appear and what evidence there is. Concepts are listed only where we can point
to the code, the schema or a test; a course topic that the project does not use is not
claimed.

\section{CE102: Problem Solving using Object Oriented Programming (Lab)}\label{sec:oop}
The backend and the portal are Java 21 programs of about 16\,000 lines (backend 11\,300, portal
4\,900), written around the object-oriented ideas of the lab course.

\begin{table}[H]
\centering\small
\caption{Object-oriented concepts in the project}\label{tab:oop}
\begin{tabularx}{\linewidth}{@{}p{3.3cm}p{5.2cm}X@{}}
\toprule \textbf{Concept} & \textbf{Where it appears} & \textbf{Why it matters here}\\ \midrule
Encapsulation & Each service hides its queries and rules behind public methods; entities have private fields with accessors; controllers never touch the database & Rules live in one place, so a rule is changed once\\
Abstraction and interfaces & \cd{NotificationProvider} interface with the implementation \cd{MockNotificationProvider} & A real SMS or e-mail provider can replace the mock without changing any caller\\
Inheritance and polymorphism & \cd{ReferenceControllers.Base<E, C>} is an abstract base class with ten subclasses for the reference resources; \cd{ApiException} extends \cd{RuntimeException} & One class provides list, get, create, update and delete for ten resources\\
Generics & \cd{Base<E, C>}, \cd{PageResponse<T>}, \cd{Db.list(Class<T>, ...)}, \cd{Db.page(...)} & Type-safe reuse without casting\\
Records & Request and response DTOs, \cd{PageParams}, \cd{PageResponse} & Immutable data carriers with built-in equality and accessors\\
Custom annotations and reflection & \cd{@Public}, \cd{@Requires}, \cd{@RequiresAny}, \cd{@ValidPassword}; \cd{AuthInterceptor} reads them at run time & Access rules are declared next to the code they protect\\
Exceptions & \cd{ApiException} with status and code; \cd{GlobalExceptionHandler}; database error translation & A runtime exception rolls the transaction back; the client gets a uniform error\\
Collections, streams, lambdas & Maps of sets for transitions; streams for filtering and mapping; lambdas in the after-commit hook & Compact, readable data handling\\
Value types & \cd{BigDecimal} for all money; \cd{java.time} for dates and times & No rounding error; clear handling of time zones (UTC)\\
Dependency injection & Constructor injection of services into services and controllers & Loose coupling and easy replacement in tests\\
Design patterns & Strategy (notification provider); Template Method (reference base class); Fluent builder (\cd{QB}); Chain of responsibility (JWT filter, then interceptor, then controller); Utility class with a private constructor (\cd{Perms}, \cd{Fees}) & Familiar structures that make the code easier to read\\
Concurrency & \cd{ConcurrentHashMap} for revoked tokens; transactions and row locks; tests that start threads together using a \cd{CountDownLatch} & Safe behaviour with many simultaneous users\\
\bottomrule
\end{tabularx}
\end{table}
\textbf{A design choice, stated plainly.} Statuses such as \texttt{ACTIVE} or
\texttt{APPROVED} are Java strings in the code and \texttt{ENUM} columns in the database,
not Java enums. This keeps the code close to the schema, and the transition tables catch
wrong values, but Java enums would add compile-time safety. We note this as a possible
improvement.

\section{CE204: Database Management}\label{sec:dbms}
This is the course the project leans on most.

\begin{table}[H]
\centering\small
\caption{Database concepts in the project}\label{tab:dbms}
\begin{tabularx}{\linewidth}{@{}p{3.4cm}p{6.1cm}X@{}}
\toprule \textbf{Concept} & \textbf{How it is used} & \textbf{Evidence}\\ \midrule
ER modelling & 12 ER diagrams covering 63 tables, generated from the live schema & Section 4.6\\
Relational schema and keys & 59 surrogate keys, 4 composite keys, 43 unique keys, 91 foreign keys & Schema dump\\
Generalisation & \cd{persons} is the parent; \cd{citizens} and \cd{employees} extend it; \cd{applicants} refer to a citizen & ER figures\\
Many-to-many & Four junction tables: user\_roles, role\_permissions, licence\_class\_assignments, challan\_violations & Table~\ref{tab:fkrules}\\
Normalisation & Third normal form with three documented exceptions & Section 4.6\\
Referential actions & 52 RESTRICT, 29 CASCADE, 10 SET NULL chosen per relationship & Table~\ref{tab:fkrules}\\
Domain integrity & 25 ENUM columns, exact \texttt{DECIMAL(10,2)} money, 26 enforced CHECK rules (7 in the base schema, 19 of the 21 in the constraints script; MySQL rejects two) & Section~\ref{sec:defects}\\
Generated columns & Three stored generated columns enforce one current row & Table~\ref{tab:gen}\\
SQL DDL and DML & 13 module scripts, the constraint, index, view, trigger and procedure scripts, and the seed scripts & \cd{database/}\\
Joins, grouping, subqueries & Inner, left and self joins; \texttt{GROUP BY ... HAVING} (queries Q1 to Q4) & Section~\ref{sec:sqltests}\\
CTEs and window functions & \texttt{WITH} queries (Q5, Q6); \texttt{RANK}, \texttt{DENSE\_RANK}, \texttt{ROW\_NUMBER}, \texttt{LAG} over partitions (Q7 to Q10) & Section~\ref{sec:sqltests}\\
Views & 6 views & Table~\ref{tab:views}\\
Indexes and query plans & 25 added indexes; \texttt{EXPLAIN} shows a scan turning into an index lookup & Section~\ref{sec:indexres}\\
Triggers & 5 triggers (history, delete guards, auto-overdue) & Table~\ref{tab:triggers}\\
Stored procedures & 3 procedures with transactions and error handlers & Table~\ref{tab:procs}\\
Transactions and ACID & Every service method is one transaction; rollback test; concurrency tests & Chapter~\ref{ch:test}\\
Locking and isolation & \cmd{SELECT ... FOR UPDATE} through \cd{Db.lock}; connections use \texttt{READ COMMITTED}; the server default is \texttt{REPEATABLE READ} & Chapter~\ref{ch:algo}\\
JSON data type & Old and new values of the audit log & Section 4.6\\
Security of the database & The application connects as a dedicated account that has rights only on the project's databases & Section~\ref{sec:envsetup}\\
\bottomrule
\end{tabularx}
\end{table}

\section{CE302: Software Engineering}\label{sec:se}
\begin{table}[H]
\centering\small
\caption{Software engineering practices in the project}\label{tab:se}
\begin{tabularx}{\linewidth}{@{}p{3.6cm}X@{}}
\toprule \textbf{Practice} & \textbf{How it was applied}\\ \midrule
Process model & Incremental development in dependency order: schema, backend, portal, database hardening, final validation (Section~\ref{sec:timeline})\\
Requirements & Numbered functional and non-functional requirements with acceptance criteria and a test that verifies each (Chapter~\ref{ch:srs})\\
Modelling (UML) & Use case, class, sequence, activity, state, component and deployment diagrams, plus data flow diagrams at three levels (Chapter~\ref{ch:design})\\
Architecture and design principles & Layered client--server design; separation of concerns; one responsibility per class; programming to an interface where a replacement is likely\\
Coding standards & Consistent package structure and naming; Bean Validation instead of ad hoc checks; conventional commit messages\\
Version control & Git with a feature branch for the portal merged into the main branch; small single-purpose commits for the database layer; secrets kept out of the repository by \texttt{.gitignore}\\
Configuration management & Environment-specific values in a git-ignored \texttt{.env} file; an example file with blank values is committed\\
Testing & Unit-level schema audit, 185 integration tests, a SQL test script, a UI crawl and a load test (Chapter~\ref{ch:test})\\
Traceability & Requirements map to tests in Table~\ref{tab:fr}; defects map to fixes in Section~\ref{sec:defects}\\
Documentation & README with set-up steps; Swagger/OpenAPI for the API; this report\\
Defect management & Defects were recorded with cause, fix and re-test (Section~\ref{sec:defects})\\
\bottomrule
\end{tabularx}
\end{table}

\section{CE202: Data Structures}
Chapter~\ref{ch:algo} gives the data structures in Table~\ref{tab:ds}. In summary:
hash maps and sets for transition tables and permissions; a sorted set to fix the order of
locks; a concurrent hash map for revoked tokens; a tree for the region hierarchy; a
directed acyclic graph for foreign-key dependencies; B$^+$-tree indexes in the database
for every key; and a table used as a durable queue for notifications. Each choice has a
stated reason and a stated cost.

\section{CE207: Design and Analysis of Algorithms}
The algorithm chapter applies the course in four ways.
\begin{itemize}
\item \textbf{Asymptotic analysis} of every critical operation (Table~\ref{tab:time}), with the assumptions stated.
\item \textbf{A correctness argument} for ordered locking: if all transactions take locks in one global order, no cycle of waiting transactions can form, so deadlock cannot occur among them.
\item \textbf{Graph algorithms.} The creation order of the 63 tables is a topological order of the foreign-key graph, computed and verified in $O(V+E)$ time; the region cycle check walks a parent chain in $O(d)$.
\item \textbf{Empirical growth.} Lookup time was measured on tables of 10\,000, 100\,000 and 500\,000 rows with and without an index (Section~\ref{sec:indexres}). The scan grows roughly in proportion to the table size and the indexed lookup stays almost flat, which agrees with $O(n)$ versus $O(\log n)$.
\end{itemize}

\section{Mapping of the Courses to the Report}
\begin{table}[H]
\centering\small
\caption{Where each course appears in the project and in this report}\label{tab:map}
\begin{tabularx}{\linewidth}{@{}p{3.3cm}p{4.6cm}X@{}}
\toprule \textbf{Course} & \textbf{Main artefacts in the project} & \textbf{Report sections}\\ \midrule
CE102 Object Oriented Programming & Backend and portal Java code; DTO records; services; annotations; interfaces & 6.1, 7.4, 7.5\\
CE204 Database Management & \cd{database/} scripts; JPA mapping; SQL tests and queries & 4.6, 5.2, 6.2, 9.7, 8.4\\
CE302 Software Engineering & SRS; UML models; Git history; test suites; defect log & 3, 4, 6.3, 7.8, 9\\
CE202 Data Structures & Maps, sets, trees, graphs, indexes, queue table & 5.4, 6.4\\
CE207 Design and Analysis of Algorithms & Locking order, topological ordering, complexity tables, index experiment & 5, 6.5, 8.4\\
\bottomrule
\end{tabularx}
\end{table}
